% Copyright (c) 2026 Eugene Joseph M. Ragasa.
% SPDX-License-Identifier: Apache-2.0
% Author: Eugene Joseph M. Ragasa
% Formal proof completed: 2026-10-20
\documentclass[11pt]{article}
\usepackage[T1]{fontenc}\usepackage{amsmath,amssymb}\usepackage[margin=1in]{geometry}\usepackage[hidelinks]{hyperref}
\title{Proof Dossier: Rank-Two Band Error Frobenius Bound}\author{Eugene Joseph M. Ragasa}\date{Formal proof completed: 2026-10-20}
\begin{document}\maketitle
\section{Identity and status}
ID \texttt{BG-FOUNDATION-THM-012}; status \texttt{proved}. Lean declaration:
\par\noindent{\tiny\path{ProjectKoios.BandGauge.Foundations.}}\\
{\tiny\path{rankTwo_max_band_error_le_frobenius}}.
\section{Source and contribution boundary}
Source-neutral supporting algebra; relation \texttt{background-only} to standard finite
Hermitian perturbation theory through pinned Mathlib. No publication-source or novelty
claim; a later diagnostic may consume this theorem without changing its authority.
\section{Statement and proof}
For two Hermitian rank-two Pauli-coordinate matrices, the maximum error between their
lower bands and between their upper bands is at most the unnormalized entrywise
Frobenius distance. The reverse triangle inequality bounds the radius change by the
Euclidean Pauli-vector change; Cauchy--Schwarz then gives the Frobenius bound.
\section{Boundary}
No factor $1/\sqrt2$ appears. The rank-normalized Frobenius loss is not thereby a Weyl
upper bound. Exact finite matrices only; no runtime evidence. Kernel axioms are \texttt{propext}, \texttt{Classical.choice}, and \texttt{Quot.sound}; no admissions. Checked with Lean \texttt{4.35.0} and Mathlib revision \texttt{56a005855ea17e0b0958f73ec801ae140b0a1679}. AI-assisted and kernel-checked. No independent expert review is recorded.
\end{document}
